\section{Resultados Esperados}

El proyecto de modernización digital de \textbf{Wayki Trek S.A.C.} se estructuró con base en metas claras, definiendo KPIs específicos y métricas agrupadas por ejes estratégicos para poder medir de forma fehaciente el éxito de la transformación digital.

\subsection{Metas Generales del Proyecto}
El plan estableció tres grandes objetivos finales concretos para evaluar el impacto global del proyecto a mediano y largo plazo:
\begin{itemize}
    \item \textbf{Optimizar la gestión de reservas en un 20\%:} Lograr una reducción sistemática en los tiempos de respuesta y una disminución de errores manuales en la atención al cliente, facilitado por el uso del orquestador \textit{Wayki Odoo CLI} y el \textit{WTK Custom Quote Wizard}.
    \item \textbf{Recuperar al menos un 10\% de carritos abandonados:} Utilizar estrategias automatizadas de reconexión (\textit{retargeting} y auto-seguimientos de Meta/WhatsApp) controladas desde el CRM para reactivar usuarios que pausaron su proceso de compra.
\end{itemize}

\subsection{Indicadores Clave de Desempeño (KPIs) por Eje Estratégico}
Para garantizar el monitoreo del avance diario y semanal, el proyecto definió la siguiente matriz de KPIs dentro del ecosistema del CRM:

\begin{itemize}
    \item \textbf{Eje de Eficiencia Operativa:}
    \begin{itemize}
        \item \textit{Indicador:} Reducción del tiempo promedio de gestión de reservas.
        \item \textit{Métrica:} Comparación del tiempo de atención y registro antes y después del despliegue.
        \item \textit{Meta:} Optimización del 20\%.
    \end{itemize}
    \item \textbf{Eje de Gestión Comercial:}
    \begin{itemize}
        \item \textit{Indicador:} Tasa de recuperación de carritos abandonados.
        \item \textit{Métrica:} Porcentaje de usuarios que completan su reserva tras recibir seguimiento.
        \item \textit{Meta:} Recuperar al menos el 10\%.
    \end{itemize}
    \item \textbf{Eje de Automatización Comercial:}
    \begin{itemize}
        \item \textit{Indicador:} Porcentaje de tareas y flujos automatizados.
        \item \textit{Métrica:} Procesos automatizados vs total de tareas manuales previas.
        \item \textit{Meta:} Automatizar al menos el 30\% de los seguimientos y comunicaciones.
    \end{itemize}
\end{itemize}

\subsection{Frecuencia de Medición e Impacto}
El esquema de auditoría de resultados y reportabilidad se estructuró con las siguientes frecuencias utilizando los Dashboards nativos de Odoo y herramientas externas (ej. Google Search Console):

\begin{itemize}
    \item \textbf{Mensual:} Incremento del número de reservas confirmadas y recuperación de carritos abandonados.
    \item \textbf{Mensual y Trimestral:} Reducción del tiempo promedio de respuesta al cliente.
    \item \textbf{Trimestral:} Tasa de conversión de leads (Win Rate).
    \item \textbf{Semestral:} Nivel de satisfacción post-reserva (encuestas automatizadas).
\end{itemize}
